\chapter{The algebraic core}
\label{chap:algebra-core}

This chapter contains the part of the proof which is independent of
schemes and cohomology.  Its inputs are finite Frobenius sets, formal
power series, sequences, finite-dimensional vector spaces, and integral
lattices.  Separating this material has two benefits.  It makes clear
which conclusions really use the geometric theorems, and it produces a
library reusable for other zeta functions and compatible systems.

\section{Formal logarithms, exponentials, and products}

Let \(R\) be a commutative \(\mathbf Q\)-algebra.  The \(T\)-adic
filtration makes \(R[[T]]\) complete and separated.  Therefore
\[
 \exp(U)=\sum_{n\geq0}\frac{U^n}{n!}
 \quad(U\in TR[[T]])
\]
and
\[
 \log(1+U)=\sum_{n\geq1}(-1)^{n+1}\frac{U^n}{n}
 \quad(U\in TR[[T]])
\]
are defined coefficientwise.  They are inverse bijections between
\(TR[[T]]\) under addition and \(1+TR[[T]]\) under multiplication.
Every coefficient depends on finitely many terms, so no analytic
convergence is involved.

An infinite product \(\prod_{\lambda}(1+u_\lambda)\), with
\(u_\lambda\in TR[[T]]\), is called \(T\)-adically locally finite if, for
each \(N\), only finitely many \(u_\lambda\) are nonzero modulo
\(T^{N+1}\).  Such a product is well defined: its image in each quotient
\(R[[T]]/(T^{N+1})\) is a finite product, and these finite products are
compatible as \(N\) varies.  Formal logarithm then converts it into the
locally finite sum of the logarithms.  This elementary completeness
argument is what justifies the closed-point Euler product.

\section{Closed points as Frobenius orbits}

Let \(Y_0\) be of finite type over \(\Fq\), and put
\(Y=Y_0\times_{\Fq}\Fbar\).  A geometric point of \(Y\) is defined over
some finite extension of \(\Fq\), so its orbit under geometric
Frobenius \(F\) is finite.  If the least positive integer fixing the
point is \(m\), the orbit has \(m\) elements and descends to a closed
point with residue field \(\mathbf F_{q^m}\).  Conversely, the geometric
fiber of a closed point \(y\) with
\(\kappa(y)=\mathbf F_{q^m}\) consists of the \(m\) embeddings of its
residue field into \(\Fbar\), and Frobenius acts transitively on them.
Thus:
\[
 \{\text{closed points of }Y_0\}
 \longleftrightarrow
 \{\text{finite Frobenius orbits in }Y(\Fbar)\},
 \tag{2.1}\label{eq:closed-orbit-bijection}
\]
with degree equal to orbit length.

Let \(a_m\) be the number of closed points of degree \(m\), and let
\(N_n=\#Y_0(\mathbf F_{q^n})\).

\begin{lemma}[Counting fixed points orbit by orbit]
\label{lem:point-count-orbits}
For every \(n\geq1\),
\[
  N_n=\sum_{m\mid n}ma_m.
  \tag{2.2}\label{eq:point-orbit-count}
\]
\end{lemma}

\begin{proof}
An \(\mathbf F_{q^n}\)-point is a geometric point fixed by \(F^n\).  On
an orbit of length \(m\), the permutation \(F^n\) is the identity if
\(m\mid n\) and has no fixed point otherwise.  In the first case all
\(m\) points contribute.  Summation over the disjoint orbits gives
\eqref{eq:point-orbit-count}.
\end{proof}

There are finitely many closed points of every bounded degree.  Indeed,
if \(m\leq r\), all geometric points above a degree-\(m\) point are
rational over \(\mathbf F_{q^{r!}}\).  They belong to the finite set
\(Y_0(\mathbf F_{q^{r!}})\).  This verifies the local finiteness needed
for the next product.

\begin{proposition}[Closed-point Euler product]
\label{prop:euler-product}
For every finite-type \(\Fq\)-scheme \(Y_0\),
\[
 \exp\left(\sum_{n\geq1}N_n\frac{T^n}{n}\right)
 =\prod_{y\in|Y_0|}(1-T^{\deg(y)})^{-1}
 \quad\text{in }1+T\mathbf Z[[T]].
 \tag{2.3}\label{eq:closed-point-product}
\]
\end{proposition}

\begin{proof}
The preceding bounded-degree argument makes the right side a
\(T\)-adically locally finite product.  Each factor has integral
coefficients, so the product belongs to \(1+T\mathbf Z[[T]]\).

Compute its logarithm in \(\mathbf Q[[T]]\).  Grouping closed points by
degree and using
\(-\log(1-U)=\sum_{s\geq1}U^s/s\), we obtain
\[
\begin{aligned}
 \log\prod_{y\in|Y_0|}(1-T^{\deg(y)})^{-1}
 &=\sum_{m\geq1}a_m\sum_{s\geq1}\frac{T^{ms}}s\\
 &=\sum_{n\geq1}\left(\sum_{m\mid n}\frac{a_m}{n/m}\right)T^n\\
 &=\sum_{n\geq1}\frac1n
        \left(\sum_{m\mid n}ma_m\right)T^n\\
 &=\sum_{n\geq1}N_n\frac{T^n}{n},
\end{aligned}
\]
where the last equality is \cref{lem:point-count-orbits}.  Formal
exponential is inverse to formal logarithm, so
\eqref{eq:closed-point-product} follows.
\end{proof}

\begin{corollary}[M\"obius recovery of closed points]
\label{cor:mobius-closed-points}
The point counts determine the number of degree-\(m\) closed points:
\[
 a_m=\frac1m\sum_{r\mid m}\mu(r)N_{m/r}.
\]
\end{corollary}

\begin{proof}
Apply ordinary M\"obius inversion to
\(N_n=\sum_{m\mid n}ma_m\).
\end{proof}

The quotient in this formula is automatically a nonnegative integer
because it counts orbits.  The integrality is much less evident if one
starts with an arbitrary numerical sequence \(N_n\); it is a geometric
constraint encoded by the Euler product.

\section{The determinant--trace identity}

Let \(K\) be a characteristic-zero field, let \(V\) be a
finite-dimensional \(K\)-vector space, and let \(u\in\End_K(V)\).

\begin{theorem}[Determinant as the exponential of traces]
\label{prop:det-trace}
In \(K[[T]]\),
\[
 \det(1-Tu)^{-1}
 =\exp\left(\sum_{n\geq1}\Tr(u^n)\frac{T^n}{n}\right).
 \tag{2.4}\label{eq:det-trace}
\]
Equivalently,
\[
 -\frac{d}{dT}\log\det(1-Tu)
 =\sum_{n\geq1}\Tr(u^n)T^{n-1}.
 \tag{2.5}\label{eq:log-derivative-trace}
\]
\end{theorem}

\begin{proof}
Scalar extension preserves determinant, trace, and equality of formal
series, so extend \(K\) to an algebraic closure.  Choose a basis in which
\(u\) is upper triangular.  Let the diagonal entries, with algebraic
multiplicity, be \(\lambda_1,\ldots,\lambda_r\).  Nilpotent entries above
the diagonal affect neither the determinant nor the traces of powers:
\[
 \det(1-Tu)=\prod_{j=1}^r(1-\lambda_jT),\qquad
 \Tr(u^n)=\sum_{j=1}^r\lambda_j^n.
\]
Thus
\[
\begin{aligned}
 -\log\det(1-Tu)
 &=-\sum_{j=1}^r\log(1-\lambda_jT)\\
 &=\sum_{j=1}^r\sum_{n\geq1}
       \frac{\lambda_j^nT^n}{n}\\
 &=\sum_{n\geq1}\Tr(u^n)\frac{T^n}{n}.
\end{aligned}
\]
Exponentiating proves \eqref{eq:det-trace}; differentiating proves
\eqref{eq:log-derivative-trace}.  Since both sides before scalar
extension have coefficients in \(K\), the identity descends to \(K\).
\end{proof}

\begin{remark}
There is also a basis-free proof.  Jacobi's determinant formula gives
\[
 \frac{d}{dT}\det(1-Tu)
 =-\det(1-Tu)\Tr((1-Tu)^{-1}u).
\]
The geometric series
\((1-Tu)^{-1}=\sum_{n\geq0}T^nu^n\) then yields
\eqref{eq:log-derivative-trace}.  The triangular proof was chosen because
it simultaneously records the reciprocal roots and makes clear that
semisimplicity is unnecessary.
\end{remark}

\begin{corollary}[Alternating determinant formula]
\label{cor:alternating-det-trace}
For a bounded graded family \((V_i,u_i)\),
\[
 \prod_i\det(1-Tu_i)^{(-1)^{i+1}}
 =\exp\left(\sum_{n\geq1}
   \left(\sum_i(-1)^i\Tr(u_i^n)\right)\frac{T^n}{n}\right).
 \tag{2.6}\label{eq:alternating-det-trace}
\]
\end{corollary}

\begin{proof}
Apply \cref{prop:det-trace} in each degree.  The factor
\(\det(1-Tu_i)^{-1}\) has logarithm
\(\sum_n\Tr(u_i^n)T^n/n\).  Raising it to \((-1)^i\) and multiplying
over \(i\) gives exactly \eqref{eq:alternating-det-trace}.  Boundedness
makes the degree product finite.
\end{proof}

\begin{lemma}[Additivity through a stable filtration]
\label{lem:det-trace-filtration}
If \(u\) preserves a finite filtration of \(V\), then
\[
 \Tr(u^n\mid V)=\sum_a\Tr(u^n\mid\operatorname{gr}_aV)
\]
and
\[
 \det(1-Tu\mid V)
 =\prod_a\det(1-Tu\mid\operatorname{gr}_aV).
\]
\end{lemma}

\begin{proof}
Choose a basis adapted to the filtration.  The matrix of \(u\) is block
upper triangular, with diagonal blocks the induced maps on the graded
pieces.  The same is true of every power.  Trace adds and determinant
multiplies across those diagonal blocks.
\end{proof}

This lemma is the linear-algebra reason spectral-sequence arguments need
only a finite Frobenius-stable filtration and subquotient control; they
do not require a noncanonical direct-sum splitting.

\section{Rational series and eventual recurrences}

Let \(a=(a_0,a_1,\ldots)\) be a sequence in a field \(K\).  Write
\[
 Sa=(a_1,a_2,\ldots)
\]
for the left shift, and let
\[
 H(a)=(a_{r+s})_{r,s\geq0}
\]
be its infinite Hankel matrix.  The rank of an infinite matrix is the
supremum of the ranks of all its finite submatrices.  Equivalently, it is
the dimension of the span of its columns when that dimension is finite.

\begin{proposition}[Rationality, recurrence, shift span, and Hankel rank]
\label{prop:rational-recurrence-hankel}
For \(f(T)=\sum_{n\geq0}a_nT^n\in K[[T]]\), the following conditions are
equivalent.
\begin{enumerate}
\item \(f\) is the Taylor expansion at zero of an element of \(K(T)\)
      regular at zero.
\item There are \(r,N\geq0\) and \(b_1,\ldots,b_r\in K\) such that
      \[
       a_n+b_1a_{n-1}+\cdots+b_ra_{n-r}=0
       \qquad(n\geq N).
       \tag{2.7}\label{eq:eventual-recurrence}
      \]
\item The \(K\)-span of \(S^na\), \(n\geq0\), is finite-dimensional.
\item The Hankel matrix \(H(a)\) has finite rank over \(K\).
\end{enumerate}
\end{proposition}

\begin{proof}
Suppose \(f=A/B\), with \(A,B\in K[T]\) and \(B(0)=1\).  Write
\[
 B(T)=1+b_1T+\cdots+b_rT^r.
\]
The coefficient of \(T^n\) in \(Bf\) is
\[
 a_n+b_1a_{n-1}+\cdots+b_ra_{n-r},
\]
where negative subscripts are interpreted as zero.  Since \(Bf=A\), this
coefficient vanishes when \(n>\deg A\).  Hence (1) implies (2).

Assume (2).  Discarding the first \(N\) coordinates, the recurrence
expresses every tail after the first \(r\) tails as a linear combination
of those \(r\).  Restoring the discarded coordinates adds only the
finite-dimensional space of sequences supported in
\(\{0,\ldots,N+r-1\}\).  Therefore all shifts lie in a
finite-dimensional span, and (2) implies (3).

The \(s\)-th column of \(H(a)\) is the sequence
\[
 (a_s,a_{s+1},a_{s+2},\ldots)=S^sa.
\]
Thus the column span of the Hankel matrix is precisely the shift span.
This proves the equivalence of (3) and (4).

Finally suppose
\[
 W=\operatorname{span}_K\{S^na:n\geq0\}
\]
is finite-dimensional.  It is stable under \(S\).  If
\[
 c(t)=t^r+c_{r-1}t^{r-1}+\cdots+c_0
\]
is the characteristic polynomial of \(S\mid W\), Cayley--Hamilton gives
\(c(S)a=0\).  In the \(n\)-th coordinate,
\[
 a_{n+r}+c_{r-1}a_{n+r-1}+\cdots+c_0a_n=0
 \qquad(n\geq0).
 \tag{2.8}\label{eq:CH-recurrence}
\]
Set
\[
 C(T)=1+c_{r-1}T+\cdots+c_0T^r.
\]
For every \(m\geq r\), the coefficient of \(T^m\) in \(Cf\) is the
left side of \eqref{eq:CH-recurrence} with \(n=m-r\), and is zero.
Therefore \(Cf\) is a polynomial of degree less than \(r\), proving that
\(f\in K(T)\).  Thus (3) implies (1).
\end{proof}

\begin{remark}
The recurrence obtained from a chosen rational expression need only be
eventual because its numerator creates finitely many exceptional
coefficients.  Passing to the full shift span absorbs those exceptions,
and Cayley--Hamilton returns a recurrence valid from \(n=0\).  This
distinction is helpful in formal code: an eventual recurrence is easy to
extract from a denominator, while a global recurrence is easy to use as
an endomorphism identity.
\end{remark}

\section{Hankel descent}

\begin{lemma}[Rank is unchanged by scalar extension]
\label{lem:matrix-rank-extension}
Let \(K\subset L\) be fields and let \(M\) be a possibly infinite matrix
with entries in \(K\).  Its rank over \(K\) equals its rank over \(L\).
\end{lemma}

\begin{proof}
For each \(r\), the assertion that the rank is at least \(r\) means that
some finite \(r\times r\) minor has nonzero determinant.  Such a
determinant lies in \(K\), and it is nonzero in \(K\) if and only if its
image is nonzero in \(L\).  The sets of possible \(r\) are therefore the
same over both fields.  Taking their supremum proves the result.
\end{proof}

\begin{theorem}[Hankel rationality descent]
\label{thm:hankel-descent}
For every field extension \(K\subset L\),
\[
  K[[T]]\cap L(T)=K(T)
\]
inside \(L[[T]]\).
\end{theorem}

\begin{proof}
Let \(f=\sum a_nT^n\) have \(a_n\in K\), and suppose
\(f\in L(T)\).  By
\cref{prop:rational-recurrence-hankel}, the Hankel matrix \(H(a)\) has
finite rank over \(L\).  Its entries lie in \(K\), so
\cref{lem:matrix-rank-extension} gives the same finite rank over \(K\).
Applying \cref{prop:rational-recurrence-hankel} again, now over \(K\),
gives \(f\in K(T)\).
\end{proof}

This argument is preferable to averaging a denominator over field
automorphisms.  A sum of denominators need not remain a denominator, and
an infinite extension may have no useful finite averaging group.  Hankel
rank is canonical and is defined by the vanishing of minors with
coefficients in the smaller field.

\section{The integral Fatou lattice}

Hankel descent gives rational coefficients in a rational expression, but
the Weil conjectures require integral polynomial factors.  The next
theorem upgrades an integral rational series to an integral numerator and
denominator.

\begin{lemma}[Finite coordinates detect a finite-dimensional sequence space]
\label{lem:finite-coordinate-detection}
Let \(V\subset\mathbf Q^{\mathbf N}\) be finite-dimensional.  There is an
\(N\) such that projection to the first \(N\) coordinates
\[
 \pi_N\colon V\longrightarrow\mathbf Q^N
\]
is injective.
\end{lemma}

\begin{proof}
The kernels \(\ker\pi_N\) form a descending sequence of subspaces of the
finite-dimensional space \(V\).  Their intersection is zero, because a
sequence all of whose coordinates vanish is the zero sequence.  The
descending sequence of dimensions eventually stabilizes.  Once it
stabilizes, the subspaces themselves stabilize, and their stable value is
their intersection.  Hence that stable value is zero.
\end{proof}

\begin{theorem}[Integral Fatou lemma]
\label{thm:integral-fatou}
If
\[
 f(T)=\sum_{n\geq0}a_nT^n\in\mathbf Z[[T]]\cap\mathbf Q(T),
\]
then \(f=A/B\) for polynomials \(A,B\in\mathbf Z[T]\) with \(B(0)=1\).
If \(a_0=1\), one also has \(A(0)=1\).
\end{theorem}

\begin{proof}
Let \(a=(a_0,a_1,\ldots)\).  Rationality and
\cref{prop:rational-recurrence-hankel} show that
\[
 V=\operatorname{span}_{\mathbf Q}\{S^na:n\geq0\}
\]
is finite-dimensional.  Choose \(N\) so that \(\pi_N\) is injective on
\(V\), using \cref{lem:finite-coordinate-detection}.

Inside \(V\), form the shift-stable abelian group
\[
 M=\sum_{n\geq0}\mathbf Z\,S^na.
 \tag{2.9}\label{eq:fatou-lattice}
\]
Every generator is an integral sequence.  Hence
\(\pi_N(M)\subset\mathbf Z^N\).  A submodule of the finitely generated
\(\mathbf Z\)-module \(\mathbf Z^N\) is finitely generated because
\(\mathbf Z\) is noetherian.  Since \(\pi_N\) is injective, \(M\) is
finitely generated as well.  It is torsion-free as a subgroup of a
\(\mathbf Q\)-vector space, so it is a finite free
\(\mathbf Z\)-module.

The shift maps \(M\) to itself.  Let
\[
 c(t)=t^r+c_{r-1}t^{r-1}+\cdots+c_0\in\mathbf Z[t]
\]
be the characteristic polynomial of the integral endomorphism
\(S\mid M\).  Cayley--Hamilton holds for endomorphisms of finite free
modules, and yields \(c(S)a=0\).  Coordinatewise,
\[
 a_{n+r}+c_{r-1}a_{n+r-1}+\cdots+c_0a_n=0
 \quad(n\geq0).
\]
Define
\[
 B(T)=1+c_{r-1}T+\cdots+c_0T^r\in\mathbf Z[T].
\]
The recurrence says that \(A(T)=B(T)f(T)\) has degree less than \(r\).
Its coefficients are integral because both \(B\) and the series \(f\)
are integral.  The constant coefficients satisfy
\[
 B(0)=1,\qquad A(0)=B(0)a_0=a_0.
\]
This proves every assertion.
\end{proof}

\begin{remark}
The essential lattice is not an arbitrarily chosen lattice in the
rational shift space.  It is the group generated by all integral shifts.
Finite-coordinate detection embeds this possibly infinite-looking group
into \(\mathbf Z^N\), where noetherianity proves finite generation.
This is the step which turns a rational recurrence into an integral monic
characteristic polynomial.
\end{remark}

\section{Normalized Gauss factors}

\begin{lemma}[Normalized rational factors are integral]
\label{lem:normalized-gauss-factor}
Suppose \(C\in\mathbf Z[T]\), \(C(0)=1\), and
\[
 C(T)=P(T)Q(T)
\]
with \(P,Q\in\mathbf Q[T]\) and \(P(0)=Q(0)=1\).  Then
\(P,Q\in\mathbf Z[T]\).
\end{lemma}

\begin{proof}
The polynomial \(C\) is primitive because its constant coefficient is
\(1\).  Write
\[
 P=aP_0,\qquad Q=bQ_0,
\]
where \(a,b\in\mathbf Q^\times\) and \(P_0,Q_0\in\mathbf Z[T]\) are
primitive.  Gauss's lemma says \(P_0Q_0\) is primitive.  Since
\[
 C=abP_0Q_0
\]
and both \(C\) and \(P_0Q_0\) are primitive, comparison of contents gives
\(ab=\pm1\), after harmlessly changing the sign of a primitive factor.
The normalizations at zero say
\[
 aP_0(0)=1,\qquad bQ_0(0)=1.
\]
Thus \(P_0(0)Q_0(0)=\pm1\).  These are integers whose product is a unit,
so each is \(\pm1\).  Hence \(a,b\in\{1,-1\}\), and \(P,Q\) are integral.
\end{proof}

\begin{corollary}[Reduced integral presentation]
\label{cor:reduced-integral-presentation}
If \(f\in1+T\mathbf Z[[T]]\cap\mathbf Q(T)\), then
\[
 f=A/B
\]
for coprime \(A,B\in\mathbf Z[T]\) satisfying
\(A(0)=B(0)=1\).
\end{corollary}

\begin{proof}
\Cref{thm:integral-fatou} gives one integral presentation with both
constant terms \(1\).  Let \(G\) be the greatest common divisor in
\(\mathbf Q[T]\), normalized by \(G(0)=1\).  Then
\[
 A=GA_1,\qquad B=GB_1,
\]
where every displayed factor has constant term \(1\).  Apply
\cref{lem:normalized-gauss-factor} to each equality.  It shows that
\(G,A_1,B_1\) are integral.  Cancel \(G\); the remaining fraction is
coprime and normalized.
\end{proof}

\section{Reciprocal determinant identities}

\begin{lemma}[Inverting an automorphism]
\label{lem:inverse-determinant}
If \(u\) is an automorphism of an \(r\)-dimensional vector space, then
\[
 \det(1-Tu)=(-T)^r\det(u)\det(1-T^{-1}u^{-1})
 \tag{2.10}\label{eq:inverse-det}
\]
in \(K[T,T^{-1}]\).
\end{lemma}

\begin{proof}
The endomorphism identity
\[
 1-Tu=(-Tu)(1-T^{-1}u^{-1})
\]
is checked by multiplication.  Taking determinants gives
\eqref{eq:inverse-det}.
\end{proof}

Let \(V,W\) have dimension \(r\), let
\(\langle\, ,\,\rangle\colon V\times W\to K\) be perfect, and suppose
automorphisms \(u,v\) satisfy
\[
 \langle ux,vw\rangle=\mu\langle x,w\rangle
 \qquad(\mu\in K^\times).
 \tag{2.11}\label{eq:similitude-pair}
\]
Identifying \(W\) with \(V^\vee\), this is equivalent to
\[
 v=\mu(u^{-1})^\vee.
 \tag{2.12}\label{eq:paired-adjoint}
\]

\begin{proposition}[Reciprocal characteristic polynomials]
\label{prop:reciprocal-characteristic}
Put
\[
 P_V(T)=\det(1-Tu\mid V),\qquad
 P_W(T)=\det(1-Tv\mid W).
\]
Then
\[
 \det(v)=\frac{\mu^r}{\det(u)}
 \tag{2.13}\label{eq:paired-determinants}
\]
and
\[
 P_W(\mu^{-1}T^{-1})
 =(-1)^rT^{-r}\det(u)^{-1}P_V(T).
 \tag{2.14}\label{eq:paired-polynomials}
\]
If the eigenvalues of \(u\) are \(\alpha_1,\ldots,\alpha_r\), those of
\(v\) are \(\mu/\alpha_1,\ldots,\mu/\alpha_r\), with multiplicity.
\end{proposition}

\begin{proof}
Equation \eqref{eq:paired-adjoint} immediately gives
\eqref{eq:paired-determinants} and the eigenvalue statement.  Over an
algebraic closure,
\[
\begin{aligned}
 P_W(\mu^{-1}T^{-1})
 &=\prod_{j=1}^r
   \left(1-\frac{\mu}{\alpha_j}\mu^{-1}T^{-1}\right)\\
 &=\prod_{j=1}^r
   \left[-(\alpha_jT)^{-1}(1-\alpha_jT)\right]\\
 &=(-1)^rT^{-r}\det(u)^{-1}P_V(T).
\end{aligned}
\]
Both sides have coefficients in \(K\), so the identity descends.
\end{proof}

\section{The graded functional-equation computation}

Let \(V_0,\ldots,V_{2d}\) be finite-dimensional spaces with
automorphisms \(F_i\).  Suppose there are perfect pairings
\[
 V_i\times V_{2d-i}\longrightarrow K
\]
of common multiplier \(\mu\), and suppose the middle pairing is
\((-1)^d\)-symmetric.  Set
\[
 b_i=\dim V_i,\quad D_i=\det(F_i),\quad
 \chi=\sum_i(-1)^ib_i,
\]
and
\[
 Z_F(T)=\prod_i\det(1-TF_i)^{(-1)^{i+1}}.
\]

\begin{theorem}[Abstract functional equation]
\label{thm:abstract-functional-equation}
One has
\[
 Z_F(\mu^{-1}T^{-1})
 =(-1)^\chi T^\chi
   \left(\prod_{i=0}^{2d}D_i^{(-1)^i}\right)Z_F(T).
 \tag{2.15}\label{eq:abstract-FE-first}
\]
If \(\mu=q^d\), then
\[
 \prod_iD_i^{(-1)^i}
 =\delta^{(-1)^d}q^{d\chi/2}
 \tag{2.16}\label{eq:alt-det-sign}
\]
for a sign \(\delta\in\{1,-1\}\).  Hence
\[
 Z_F(q^{-d}T^{-1})
 =\varepsilon q^{d\chi/2}T^\chi Z_F(T),
 \quad
 \varepsilon=(-1)^\chi\delta^{(-1)^d}.
 \tag{2.17}\label{eq:abstract-FE-final}
\]
\end{theorem}

\begin{proof}
Apply \cref{prop:reciprocal-characteristic} to the pair
\((V_i,V_{2d-i})\).  Since \(i\) and \(2d-i\) have the same parity,
\[
\begin{aligned}
 Z_F(\mu^{-1}T^{-1})
 &=\prod_i
 \det(1-\mu^{-1}T^{-1}F_{2d-i})^{(-1)^{i+1}}\\
 &=\prod_i
 \left((-1)^{b_i}T^{-b_i}D_i^{-1}
       \det(1-TF_i)\right)^{(-1)^{i+1}}.
\end{aligned}
\]
The last determinant factors give \(Z_F(T)\).  The exponent of \(T\) is
\[
 -\sum_i(-1)^{i+1}b_i=\chi.
\]
The product of the signs is \((-1)^\chi\): a sign raised to exponent
\(-1\) is the same sign.  The remaining determinants contribute
\(\prod_iD_i^{(-1)^i}\).  This proves
\eqref{eq:abstract-FE-first}.

For \(i<d\), \eqref{eq:paired-determinants} says
\[
 D_iD_{2d-i}=\mu^{b_i}.
 \tag{2.18}\label{eq:det-pairs}
\]
In degree \(d\), taking determinants in the middle similitude relation
gives
\[
 D_d^2=\mu^{b_d}.
 \tag{2.19}\label{eq:middle-det-square}
\]
If \(d\) is odd, the middle pairing is alternating, so \(b_d\) is even.
If \(d\) is even, \(db_d\) is even without a restriction on \(b_d\).
Thus \(db_d/2\) is always an integer.  With \(\mu=q^d\), put
\[
 \delta=D_dq^{-db_d/2}.
\]
Equation \eqref{eq:middle-det-square} gives \(\delta^2=1\).

Pairing the determinant terms symmetrically about degree \(d\) now gives
\[
\begin{aligned}
 \prod_iD_i^{(-1)^i}
 &=\prod_{i<d}(D_iD_{2d-i})^{(-1)^i}
     D_d^{(-1)^d}\\
 &=\delta^{(-1)^d}
   q^{d\left(\sum_{i<d}(-1)^ib_i+
           \frac12(-1)^db_d\right)}.
\end{aligned}
\]
Dimension symmetry \(b_i=b_{2d-i}\) gives
\[
 \frac\chi2
 =\sum_{i<d}(-1)^ib_i+\frac12(-1)^db_d.
\]
Substitution proves \eqref{eq:alt-det-sign}, and then
\eqref{eq:abstract-FE-final}.
\end{proof}

\begin{remark}
The notation \(q^{d\chi/2}\) hides no fractional power.  If \(d\) is
even, \(d\chi\) is even.  If \(d\) is odd, the middle alternating form
forces \(b_d\), and hence \(\chi\), to be even.  The proof above
constructs the exponent as an integer before using exponential notation.
\end{remark}

\begin{remark}
For an alternating middle pairing, the Pfaffian identity improves
\(D_d^2=\mu^{b_d}\) to \(D_d=\mu^{b_d/2}\), so \(\delta=1\).  For a
symmetric middle pairing both signs can occur.  The Weil conjectures need
only the uniform conclusion \(\varepsilon\in\{1,-1\}\).
\end{remark}

\section{Formalization interfaces}

\begin{formalization}
The algebraic development should be divided into small, geometry-free
modules.
\begin{enumerate}
\item A \(T\)-adically locally finite product API should prove that
      formal logarithm exchanges Euler products and divisor sums.
\item A finite-permutation API should identify fixed points of \(F^n\)
      with orbit lengths dividing \(n\), yielding
      \cref{prop:euler-product}.
\item Determinants and traces should provide
      \cref{prop:det-trace}, its graded form, and filtration
      multiplicativity.
\item A sequence module should identify rational series, eventual
      recurrences, finite shift spans, and finite Hankel rank.  Matrix
      rank under scalar extension then gives
      \cref{thm:hankel-descent}.
\item The Fatou module should construct the integral shift lattice
      \eqref{eq:fatou-lattice}, prove finite generation by finite-coordinate
      detection, and apply Cayley--Hamilton over \(\mathbf Z\).
\item A polynomial-content module should expose the normalized factor
      lemma and reduced integral presentation separately.
\item A `PerfectSimilitudePairing` should derive adjoint, determinant,
      eigenvalue, and reciprocal-polynomial identities from one
      compatibility equation.
\item A bounded graded pairing should return
      \cref{thm:abstract-functional-equation}, including the middle-form
      parity lemma and an integer exponent for \(q\).
\end{enumerate}
These results should import no scheme theory.  The cohomological chapters
instantiate them with closed points, Frobenius modules, trace formula,
and Poincare duality.
\end{formalization}
